\documentclass[aspectratio=169, xcolor=dvipsnames]{beamer}

% --- Theme Setup ---
\usetheme{Madrid}
\usecolortheme{default}

% --- Packages ---
\usepackage{graphicx}
\usepackage{xcolor}
\usepackage{booktabs}
\usepackage{hyperref}
\usepackage{adjustbox}
\usepackage{amssymb}
\usepackage{amsmath}
\usepackage{listings}
\usepackage{caption}
\usepackage{tikz}

\lstset{
  basicstyle=\ttfamily\small,
  keywordstyle=\color{blue}\bfseries,
  stringstyle=\color{red},
  showstringspaces=false,
  breaklines=true
}

% --- Title Information ---
\title{Module 43}
\subtitle{Indexing and Hashing/3: Indexing/3}
\author{Partha Pratim Das}
\institute{Department of Computer Science and Engineering\\Indian Institute of Technology, Kharagpur\\ppd@cse.iitkgp.ac.in}
\date{\today}

\begin{document}

% Title Slide
\begin{frame}
    \titlepage
\end{frame}

\begin{frame}{Module Recap}
    \begin{block}{}
        \begin{itemize}[<+->]
            \item Recapitulated the notions of Balanced Binary Search Trees as options for optimal in-memory search data structures
            \item Understood the issues relating to external data structures for persistent data
            \item Explored 2-3-4 Tree in depth as a precursor to B/B+-Tree for an efficient external data structure for database and index tables
        \end{itemize}
    \end{block}
\end{frame}

\begin{frame}{Module Objectives}
    \begin{block}{}
        \begin{itemize}[<+->]
            \item To understand the design of B+ Tree Index Files as a generalization of 2-3-4 Tree
            \item To understand the fundamentals of B-Tree Index Files
        \end{itemize}
    \end{block}
\end{frame}

\begin{frame}{Module Outline}
    \begin{block}{}
        \begin{itemize}[<+->]
            \item B+ Tree Index Files
            \item B-Tree Index Files
        \end{itemize}
    \end{block}
\end{frame}

\section{B+ Tree Index Files}
\begin{frame}
  \vfill
  \centering
  \begin{beamercolorbox}[sep=8pt,center,shadow=true,rounded=true]{title}
    \usebeamerfont{title}B+ Tree Index Files\par%
  \end{beamercolorbox}
  \vfill
\end{frame}

\begin{frame}{The B+ Tree}
    \begin{itemize}[<+->]
        \item Is a balanced binary search tree
        \item Follows a multi-level index format like 2-3-4 Tree
        \item Has the leaf nodes denoting actual data pointers
        \item Ensures that all leaf nodes remain at the same height (like 2-3-4 Tree)
        \item Has the leaf nodes linked using a link list
        \item Can support random access as well as sequential access
    \end{itemize}
\end{frame}

\begin{frame}{The B+ Tree (2)}
    \begin{itemize}[<+->]
        \item \textbf{Internal node contains}
        \begin{itemize}[<+->]
            \item At least $\lceil n/2 \rceil$ child pointers, except the root node
            \item At most $n$ pointers
        \end{itemize}
        \item \textbf{Leaf node contains}
        \begin{itemize}[<+->]
            \item At least $\lceil n/2 \rceil$ record pointers and $\lceil n/2 \rceil$ key values
            \item At most $n$ record pointer and $n$ key values
            \item One block pointer $P$ to point to next leaf node
        \end{itemize}
    \end{itemize}
\end{frame}

\begin{frame}{B+ Tree (3): Search}
    \begin{itemize}[<+->]
        \item Suppose we have to search 55 in the B+ tree below
        \item First, we will fetch for the intermediary node which will direct to the leaf node that can contain a record for 55
        \item So, in the intermediary node, we will find a branch between 50 and 75 nodes
        \item Then at the end, we will be redirected to the third leaf node
        \item Here DBMS will perform a sequential search to find 55
    \end{itemize}
\end{frame}

\begin{frame}{B+ Tree: Insert}
    \begin{itemize}[<+->]
        \item Suppose we want to insert a record 60 that goes to 3rd leaf node after 55
        \item The leaf node of this tree is already full, so we cannot insert 60 there
        \item So we have to split the leaf node, so that it can be inserted into tree without affecting the fill factor, balance and order
        \item The 3rd leaf node has the values (50, 55, 60, 65, 70) and its current root node is 50
        \item We will split the leaf node of the tree in the middle so that its balance is not altered
    \end{itemize}
\end{frame}

\begin{frame}{B+ Tree: Insert (2)}
    \begin{itemize}[<+->]
        \item So we can group (50, 55) and (60, 65, 70) into 2 leaf nodes
        \item If these two has to be leaf nodes, the intermediate node cannot branch from 50
        \item It should have 60 added to it, and then we can have pointers to a new leaf node
        \item This is how we can insert an entry when there is overflow. In a normal scenario, it is very easy to find the node where it fits and then place it in that leaf node
    \end{itemize}
\end{frame}

\begin{frame}{B+ Tree: Delete}
    \begin{itemize}[<+->]
        \item To delete 60, we have to remove 60 from intermediate node as well as 4th leaf node
        \item If we remove it from the intermediate node, then the tree will not remain a B+ tree
        \item So with deleting 60 we re-arranging the nodes
    \end{itemize}
\end{frame}

\begin{frame}{B+ Tree Index Files: Advantages \& Disadvantages}
    \begin{alertblock}{Important}
        \begin{itemize}[<+->]
            \item B+ tree indices are an alternative to indexed-sequential files
            \item \textbf{Disadvantage of ISAM files}: Performance degrades as file grows, since many overflow blocks get created. Periodic reorganization of entire file is required
            \item \textbf{Advantage of B+ tree index files}: Automatically reorganizes itself with small, local, changes, in the face of insertions and deletions. Reorganization of entire file is not required to maintain performance
            \item \textbf{(Minor) disadvantage of B+ trees}: Extra insertion and deletion overhead, space overhead
            \item Advantages of B+ trees outweigh disadvantages; B+ trees are used extensively
        \end{itemize}
    \end{alertblock}
\end{frame}

\begin{frame}{B+ Tree Index Files: Structure}
    \begin{itemize}[<+->]
        \item All paths from root to leaf are of the same length
        \item Each node that is not a root or a leaf has between $\lceil n/2 \rceil$ and $n$ children
        \item A leaf node has between $\lceil (n-1)/2 \rceil$ and $n-1$ values
        \item \textbf{Special cases}:
        \begin{itemize}[<+->]
            \item If the root is not a leaf, it has at least 2 children.
            \item If the root is a leaf (that is, there are no other nodes in the tree), it can have between 0 and $(n-1)$ values.
        \end{itemize}
    \end{itemize}
\end{frame}

\begin{frame}{B+ Tree Index Files: Node Structure}
    \begin{itemize}[<+->]
        \item $K_i$ are the search-key values
        \item $P_i$ are pointers to children (for non-leaf nodes) or pointers to records or buckets of records (for leaf nodes).
        \item The search-keys in a node are ordered
        \item $K_1 < K_2 < K_3 < \dots < K_{n-1}$ (Initially assume no duplicate keys, address duplicates later)
    \end{itemize}
    \begin{center}
        $P_1, K_1, P_2, \dots, P_{n-1}, K_{n-1}, P_n$
    \end{center}
\end{frame}

\begin{frame}{B+ Tree Index Files: Leaf Nodes}
    \begin{itemize}[<+->]
        \item For $i = 1, 2, \dots, n-1$, pointer $P_i$ points to a file record with search-key value $K_i$
        \item If $L_i, L_j$ are leaf nodes and $i < j$, $L_i$'s search-key values are less than or equal to $L_j$'s search-key values
        \item $P_n$ points to next leaf node in search-key order
    \end{itemize}
\end{frame}

\begin{frame}{B+ Tree Index Files: Non-Leaf Nodes}
    \begin{itemize}[<+->]
        \item Non leaf nodes form a multi-level sparse index on the leaf nodes. For a non-leaf node with $m$ pointers:
        \item All the search-keys in the subtree to which $P_1$ points are less than $K_1$
        \item For $2 \le i \le n-1$, all the search-keys in the subtree to which $P_i$ points have values greater than or equal to $K_{i-1}$ and less than $K_i$
        \item All the search-keys in the subtree to which $P_n$ points have values greater than or equal to $K_{n-1}$
    \end{itemize}
\end{frame}

\begin{frame}{B+ Tree Index Files: Observations}
    \begin{itemize}[<+->]
        \item Since the inter-node connections are done by pointers, logically close blocks need not be physically close
        \item The non-leaf levels of the B+ tree form a hierarchy of sparse indices
        \item The B+ tree contains a relatively small number of levels
        \item Level below root has at least $2 \times \lceil n \rceil$ values
        \item Next level has at least $2 \times \lceil n \rceil \times \lceil n/2 \rceil$ values, etc.
        \item If there are $K$ search-key values in the file, the tree height is no more than $\lceil \log_{\lceil n/2 \rceil} (K) \rceil$
        \item thus searches can be conducted efficiently
        \item Insertions and deletions to the main file can be handled efficiently, as the index can be restructured in logarithmic time
    \end{itemize}
\end{frame}

\begin{frame}{B+ Tree Index Files: Queries}
    \begin{itemize}[<+->]
        \item \textbf{Find record with search-key value V}
        \item \textbf{a)} $C$ = root
        \item \textbf{b)} While $C$ is not a leaf node
        \begin{itemize}
            \item i) Let $i$ be least value such that $V \le K_i$
            \item ii) If no such exists, set $C$ = last non-null pointer in $C$
            \item iii) Else \{ if ($V = K_i$) Set $C = P_{i+1}$ else set $C = P_i$ \}
        \end{itemize}
        \item \textbf{c)} Let $i$ be least value s.t. $K_i = V$
        \item \textbf{d)} If there is such a value $i$, follow pointer $P_i$ to the desired record
        \item \textbf{e)} Else no record with search-key value $k$ exists
    \end{itemize}
\end{frame}

\begin{frame}{B+ Tree Index Files: Queries (2)}
    \begin{itemize}[<+->]
        \item If there are $K$ search-key values in the file, the height of the tree is no more than $\lceil \log_{\lceil n/2 \rceil} (K) \rceil$
        \item A node is generally the same size as a disk block, typically 4 kilobytes
        \item and $n$ is typically around 100 (40 bytes per index entry)
        \item With 1 million search key values and $n = 100$: at most $\log_{50}(1,000,000) = 4$ nodes are accessed in a lookup
        \item Contrast this with a balanced binary tree with 1 million search key values - around 20 nodes are accessed in a lookup
        \item above difference is significant since every node access may need a disk I/O, costing around 20 milliseconds
    \end{itemize}
\end{frame}

\begin{frame}{B+ Tree Index Files: Handling Duplicates}
    \begin{itemize}[<+->]
        \item With duplicate search keys
        \item In both leaf and internal nodes: we cannot guarantee that $K_1 < K_2 < \dots < K_{n-1}$, but can guarantee $K_1 \le K_2 \le \dots \le K_{n-1}$
        \item Search-keys in the subtree to which $P_i$ points
        \begin{itemize}
            \item $\triangleright$ are $\le K_i$, but not necessarily $< K_i$
            \item $\triangleright$ To see why, suppose same search key value $V$ is present in two leaf node $L_i$ and $L_{i+1}$. Then in parent node $K_i$ must be equal to $V$
        \end{itemize}
    \end{itemize}
\end{frame}

\begin{frame}{Updates on B+ Trees: Insertion}
    \begin{itemize}[<+->]
        \item Find the leaf node in which the search-key value would appear
        \item If the search-key value is already present in the leaf node
        \begin{itemize}
            \item Add record to the file
            \item If necessary add a pointer to the bucket
        \end{itemize}
        \item If the search-key value is not present, then
        \begin{itemize}
            \item Add the record to the main file (and create a bucket if necessary)
            \item If there is room in the leaf node, insert (key-value, pointer) pair in the leaf node
            \item Otherwise, split the node (along with the new (key-value, pointer) entry)
        \end{itemize}
    \end{itemize}
\end{frame}

\begin{frame}{Updates on B+ Trees: Insertion (2)}
    \begin{itemize}[<+->]
        \item \textbf{Splitting a leaf node}:
        \item take the $n$ (search-key value, pointer) pairs (including the one being inserted) in sorted order. Place the first $\lceil n/2 \rceil$ in the original node, and the rest in a new node
        \item let the new node be $p$, and let $k$ be the least key value in $p$. Insert $(k, p)$ in the parent of the node being split
        \item If the parent is full, split it and propagate the split further up
        \item Splitting of nodes proceeds upwards till a node that is not full is found
        \item In the worst case the root node may be split increasing the height of the tree by 1
    \end{itemize}
\end{frame}

\begin{frame}{Updates on B+ Trees: Deletion}
    \begin{itemize}[<+->]
        \item Find the record to be deleted, and remove it from the main file and from the bucket
        \item Remove (search-key value, pointer) from the leaf node if there is no bucket or if the bucket has become empty
        \item If the node has too few entries due to the removal, and the entries in the node and a sibling fit into a single node, then \textbf{merge siblings}:
        \begin{itemize}
            \item Insert all the search-key values in the two nodes into a single node (the one on the left), and delete the other node.
            \item Delete the pair $(K_{i-1}, P_i)$, where $P_i$ is the pointer to the deleted node, from its parent, recursively.
        \end{itemize}
    \end{itemize}
\end{frame}

\begin{frame}{Updates on B+ Trees: Deletion (2)}
    \begin{itemize}[<+->]
        \item Otherwise, if the node has too few entries due to the removal, but the entries in the node and a sibling do not fit into a single node, then \textbf{redistribute pointers}:
        \begin{itemize}
            \item Redistribute the pointers between the node and a sibling such that both have more than the minimum number of entries
            \item Update the corresponding search-key value in the parent of the node
        \end{itemize}
        \item The node deletions may cascade upwards till a node which has $\lceil n/2 \rceil$ or more pointers is found
        \item If the root node has only one pointer after deletion, it is deleted and the sole child becomes the root
    \end{itemize}
\end{frame}

\begin{frame}{B+ Tree File Organization}
    \begin{itemize}[<+->]
        \item Index file degradation problem is solved by using B+ Tree indices
        \item Data file degradation problem is solved by using B+ Tree File Organization
        \item The leaf nodes in a B+ tree file organization store records, instead of pointers
        \item Leaf nodes are still required to be half full
        \item Since records are larger than pointers, the maximum number of records that can be stored in a leaf node is less than the number of pointers in a non-leaf node
        \item Insertion and deletion are handled in the same way as insertion and deletion of entries in a B+ tree index
    \end{itemize}
\end{frame}

\begin{frame}{Non-Unique Search Keys \& Relocation}
    \begin{itemize}[<+->]
        \item Alternatives to scheme described earlier
        \begin{itemize}
            \item Buckets on separate block (bad idea)
            \item List of tuple pointers with each key
            \item Make search key unique by adding a record-identifier (widely used)
        \end{itemize}
        \item \textbf{Record Relocation and Secondary Indices}
        \begin{itemize}
            \item If a record moves, all secondary indices that store record pointers have to be updated
            \item Solution: Use primary-index search key instead of record pointer in secondary index
            \item Extra traversal of primary index to locate record, but node splits are cheap
        \end{itemize}
    \end{itemize}
\end{frame}

\begin{frame}{Indexing Strings}
    \begin{itemize}[<+->]
        \item Variable length strings as keys
        \item Variable fanout
        \item Use space utilization as criterion for splitting, not number of pointers
        \item \textbf{Prefix compression}
        \begin{itemize}[<+->]
            \item Key values at internal nodes can be prefixes of full key
            \item $\triangleright$ Keep enough characters to distinguish entries in the subtrees separated by the key value
            \item - For example, 'Silas' and 'Silberschatz' can be separated by 'Silb'
            \item Keys in leaf node can be compressed by sharing common prefixes
        \end{itemize}
    \end{itemize}
\end{frame}

\section{B-Tree Index Files}
\begin{frame}
  \vfill
  \centering
  \begin{beamercolorbox}[sep=8pt,center,shadow=true,rounded=true]{title}
    \usebeamerfont{title}B-Tree Index Files\par%
  \end{beamercolorbox}
  \vfill
\end{frame}

\begin{frame}{B-Tree Index Files}
    \begin{itemize}[<+->]
        \item Similar to B+ tree, but B-tree allows search-key values to appear only once; eliminates redundant storage of search keys
        \item Search keys in non-leaf nodes appear nowhere else in the B-tree; an additional pointer field for each search key in a non-leaf node must be included
    \end{itemize}
\end{frame}

\begin{frame}{Module Summary}
    \begin{block}{}
        \begin{itemize}[<+->]
            \item Understood the design of B+ Tree Index Files as a generalization of 2-3-4 Tree
            \item Investigated the fundamentals of B-Tree Index Files
        \end{itemize}
    \end{block}
    \vspace{2em}
    \begin{center}
    \small \textit{Slides used in this presentation are borrowed from \url{http://db-book.com/} with kind permission of the authors.} \\
    \small \textit{Edited and new slides are marked with 'PPD'.}
    \end{center}
\end{frame}

\end{document}
